\section*{Bab \ref{ch:b2:conjugate-additions} --- Karbonil Terkonjugasi: Michael dan Robinson}

\begin{solution}{exo:b2:conjugate-additions:1}
\begin{itemize}
  \item \ce{CH3MgBr}: terutama 1,2 (keras, ireversibel), 1-metilsikloheks-2-en-1-ol.
  \item \ce{(CH3)2CuLi}: 1,4, 3-metilsikloheksanon.
  \item Tiofenol dengan basa: 1,4 (tiolat bersifat lunak), 3-(fenilsulfanil)sikloheksanon.
  \item \ce{LiAlH4}: 1,2, sikloheks-2-en-1-ol.
\end{itemize}
\end{solution}

\begin{solution}{exo:b2:conjugate-additions:2}
Donor: pentana-2,4-dion, nitrometana, dietil propanadioat (\ce{C-H} yang
asam). Akseptor: propenanitril, metil propenoat, but-3-en-2-on (\ce{C=C}
yang terkonjugasi dengan gugus akseptor).
\end{solution}

\begin{solution}{exo:b2:conjugate-additions:3}
\ce{C4H9Br + 2Li -> C4H9Li + LiBr}, lalu \ce{2C4H9Li + CuI -> (C4H9)2CuLi + LiI}, dalam eter pada
suhu rendah di bawah gas inert.
\end{solution}

\begin{solution}{exo:b2:conjugate-additions:4}
Dietil (3-oksobutil)propanadioat, \ce{(EtOOC)2CHCH2CH2COCH3}. Hidrolisis
menghasilkan asam propanadioat tersubstitusi, yang melepaskan \ce{CO2} pada
pemanasan (\cref{prop:b2:enolates-aldol:decarboxylation}):
asam 5-oksoheksanoat, \ce{HOOCCH2CH2CH2COCH3}.
\end{solution}

\begin{solution}{exo:b2:conjugate-additions:5}
Litium dimetilkuprat dalam eter pada suhu rendah, lalu pengolahan berair:
\omterm{def:b2:conjugate-additions:addition-modes}{adisi konjugat}. Metilmagnesium bromida, yang lebih keras, terutama beradisi
1,2 dan menghasilkan 1-metilsikloheks-2-en-1-ol sebagai produk utama.
\end{solution}

\begin{solution}{exo:b2:conjugate-additions:6}
\ce{Et3N} melepaskan proton dari sedikit \ce{EtSH} sehingga terbentuk
\ce{EtS-}; tiolat itu beradisi pada karbon $\beta$, \omterm{def:b2:enolates-aldol:enolate}{enolat} yang terbentuk
mengambil proton dari \ce{EtSH} lain, yang meregenerasi \ce{EtS-}. Produk:
4-(etilsulfanil)butan-2-on. Belerang berukuran besar dan mudah
terpolarisasi, dengan pasangan elektron bebas berenergi tinggi: sebuah
\omterm{def:b2:frontier-orbitals:hard-soft}{nukleofil lunak}, di bawah \omterm{def:b2:frontier-orbitals:control}{kendali orbital}, sehingga menyerang di tempat \omterm{def:b2:frontier-orbitals:frontier}{LUMO}
paling besar, dan adisinya reversibel, yang juga mendukung adduk-1,4 yang
lebih stabil.
\end{solution}

\begin{solution}{exo:b2:conjugate-additions:7}
Sianida menyerang karbon karbonil yang lebih positif dengan lebih cepat:
sianohidrin adalah produk kinetik. Pembentukannya reversibel; pada suhu yang
lebih tinggi, jika diberi waktu, sianida dilepaskan dan beradisi kembali
pada karbon $\beta$, sehingga terbentuk adduk yang mempertahankan ikatan $\pi$
\ce{C=O} yang lebih kuat, yaitu produk termodinamik.
\end{solution}

\begin{solution}{exo:b2:conjugate-additions:8}
Dengan 2-metilsikloheksana-1,3-dion: keton Wieland--Miescher (gambar \omterm{def:b2:conjugate-additions:robinson}{anulasi
Robinson}). Dengan sikloheksanon: 4,4a,5,6,7,8-heksahidronaftalen-2(3\textit{H})-on,
sebuah enon bisiklik yang \ce{C=C}-nya berakhir pada karbon fusi cincin,
tanpa gugus metil angular.
\end{solution}

\begin{solution}{exo:b2:conjugate-additions:9}
\ce{CH3COCH2CH2CH2COCH3}: C2 dan C6 adalah karbonilnya, lima karbon dari C2
sampai C6 jika keduanya ikut dihitung, sebuah 1,5-diketon. Potong C3--C4:
donornya \omterm{def:b2:enolates-aldol:enolate}{enolat} propanon pada C3 (lebih baik etil 3-oksobutanoat, yang
esternya disingkirkan kemudian melalui hidrolisis dan \omterm{def:b2:enolates-aldol:decarboxylation}{dekarboksilasi});
akseptornya but-3-en-2-on. Kondisi: etil 3-oksobutanoat, but-3-en-2-on,
natrium etoksida katalitik dalam etanol; lalu asam berair dan pemanasan.
\end{solution}

\begin{solution}{exo:b2:conjugate-additions:10}
\omterm{def:b2:enolates-aldol:aldol}{Aldol} membentuk ikatan \ce{C-C} tetapi mengubah ikatan $\pi$ \ce{C=O} menjadi
ikatan $\sigma$ \ce{C-O}; neraca energinya kecil dan dua molekul menjadi
satu, sehingga $K$ sedang saja dan \omterm{def:b2:enolates-aldol:aldol}{retro-aldol} mudah terjadi. \omterm{def:b2:conjugate-additions:michael}{Adisi Michael}
membentuk ikatan \ce{C-C} dengan mengorbankan ikatan $\pi$ \ce{C=C} yang
lebih lemah dan mempertahankan kedua
\ce{C=O}: adisi itu jelas menguntungkan. Reaksi baliknya akan memerlukan
\omterm{def:b2:enolates-aldol:enolate}{enolat} keton untuk mengusir anion propanadioat yang terstabilkan;
kesetimbangannya terletak begitu jauh di sisi adduk sehingga dengan basa
katalitik adduk itu bertahan.
\end{solution}

\begin{solution}{exo:b2:conjugate-additions:11}
\ce{C(COOEt)2(CH2CH2COOCH3)2}. Setelah \omterm{def:b2:conjugate-additions:michael}{adisi Michael} pertama, karbon
pusatnya masih mempunyai satu hidrogen di antara dua ester, yang sama
asamnya seperti sebelumnya: basa melepaskannya dan adisi kedua menyusul.
\end{solution}

\begin{solution}{exo:b2:conjugate-additions:12}
Melalui \omterm{def:b2:enolates-aldol:enolate}{enolat} yang lebih tersubstitusi (ke arah karbon yang membawa metil):
oktalon dengan gugus metil pada karbon fusi cincin,
4a-metil-4,4a,5,6,7,8-heksahidronaftalen-2(3\textit{H})-on. Melalui \omterm{def:b2:enolates-aldol:enolate}{enolat}
yang kurang tersubstitusi (pada sisi \ce{CH2}): gugus metil berakhir pada
karbon cincin di sebelah fusi. Produk pertama lebih disukai dalam kondisi
yang menyeimbangkan (alkoksida dalam alkoholnya, atau jalur \omterm{def:b2:amines:imine}{enamina} dengan
\omterm{def:b2:amines:class}{amina sekunder} dan asam), yang menghasilkan \omterm{def:b2:enolates-aldol:kinetic-enolate}{enolat termodinamik} yang lebih
tersubstitusi (\cref{prop:b2:enolates-aldol:enolate-control}).
\end{solution}

\begin{solution}{pb:b2:conjugate-additions:1}
\textbf{1.} Cincin sikloheksana dengan \ce{C=O} pada C1 dan C3 dan gugus
metil pada C2; hidrogen pada C2 adalah yang paling asam.
\textbf{2.} Pelepasannya menghasilkan anion yang muatannya terdelokalisasi
pada C2 dan kedua oksigen; pada sikloheksanon hanya satu oksigen yang ikut
memikulnya, dan keasamannya terlalu lemah untuk diukur dalam air.
\textbf{3.} Muatan pada C2; \ce{C1=C2} dengan \ce{O^-} pada C1; \ce{C2=C3} dengan \ce{O^-} pada C3.
\textbf{4.} Ya: dengan hidroksida, $K = 10^{14.0 - 5.3} = 10^{8.7}$ (air
sebagai asam konjugasi, dari
$K_e$); dengan etoksida, $10^{15.9-5.3} = 10^{10.6}$. Keduanya tuntas.
% ledger: ke:25C, pka:ethanol
\textbf{5.} \ce{C1(OH)=C2(CH3)}: C2 masih membawa satu hidrogen, dan hanya
itulah yang diperlukan sebuah \omterm{def:b2:enolates-aldol:tautomerism}{enol}; gugus metil hanya menggantikan hidrogen
kedua.
\textbf{6.} Lunak: muatannya tersebar pada tiga atom dan anionnya
terstabilkan. \omterm{def:b2:enolates-aldol:enolate}{Enolat} itu menyerang karbon
$\beta$ (\cref{prop:b2:conjugate-additions:regio}).
\textbf{7.} C2 \omterm{def:b2:enolates-aldol:enolate}{enolat} berikatan dengan ujung \ce{CH2} but-3-en-2-on; \omterm{def:b2:enolates-aldol:enolate}{enolat}
yang terbentuk mengambil sebuah proton:
2-metil-2-(3-oksobutil)sikloheksana-1,3-dion.
\textbf{8.} Setiap \omterm{def:b2:enolates-aldol:enolate}{enolat} baru mengambil proton dari sebuah molekul dion
(atau dari pelarut) dan dengan demikian meregenerasi \omterm{def:b2:enolates-aldol:enolate}{enolat} donor: basanya
tidak terpakai.
\textbf{9.} \omterm{def:b2:enolates-aldol:enolate}{Enolat} itu ambiden; dengan elektrofil karbon yang lunak, di bawah
\omterm{def:b2:frontier-orbitals:control}{kendali orbital}, \omterm{def:b2:enolates-aldol:enolate}{enolat} itu bereaksi pada karbon, dan adduk \ce{C-C}
mempertahankan dua ikatan \ce{C=O}, kombinasi yang lebih kuat; adduk
\ce{O}, seandainya terbentuk, akan berbalik.
\textbf{10.} Satu: C2, yang terikat pada C1, C3, metil, dan rantai.
\textbf{11.} Karbon karbonil rantai dan karbonil cincin C1 (atau C3): \ce{C(=O)}--\ce{CH2}--\ce{CH2}--C2--C1,
lima karbon jika ujung-ujungnya ikut dihitung.
\textbf{12.} Untuk menghabiskan seluruh dion, pereaksi yang lebih berharga;
hanya sedikit, karena but-3-en-2-on sangat beracun dan berpolimerisasi.
\textbf{13.} C4 dan C6 cincin (di sebelah C3 dan C1), \ce{CH2} rantai di
sebelah karbonilnya, dan \ce{CH3} ujung rantai. C2 tidak mempunyai satu pun.
\textbf{14.} \omterm{def:b2:enolates-aldol:enolate}{Enolat} \ce{CH3} ujung menyerang sebuah karbonil cincin: cincin
yang tertutup terdiri atas karbon itu, karbon karbonil rantai, dua \ce{CH2},
C2, dan karbon karbonil yang diserang: enam atom.
\textbf{15.} \omterm{def:b2:enolates-aldol:enolate}{Enolat} \ce{CH2} bagian dalam rantai akan menutup cincin
beranggota empat, yang tegang. \omterm{def:b2:enolates-aldol:enolate}{Enolat} cincin (C4 atau C6) yang menyerang
karbonil rantai juga menutup cincin beranggota enam, tetapi berjembatan:
\omterm{def:b2:enolates-aldol:aldol}{aldol} itu tidak dapat terdehidrasi (ikatan rangkapnya akan terletak pada
kepala jembatan sistem bisiklik yang kecil) dan berbalik, sedangkan \omterm{def:b2:enolates-aldol:aldol}{aldol}
yang terfusi terus dikuras oleh dehidrasi.
\textbf{16.} \omterm{def:b2:enolates-aldol:aldol}{Aldol} itu membawa \ce{OH} pada bekas karbon karbonil cincin,
yang kini berada pada fusi cincin; lepasnya air bersama sebuah hidrogen dari
bekas karbon \ce{CH3} (kini $\alpha$ terhadap keton rantai) menghasilkan
\ce{C=C} yang terkonjugasi dengan keton itu.
\textbf{17.} Ikatan rangkap yang baru terkonjugasi dengan karbonil: eliminasi
E1cB melalui \omterm{def:b2:enolates-aldol:enolate}{enolat}, yang didukung oleh pemanasan
(\cref{prop:b2:enolates-aldol:aldol-equilibrium}).
\textbf{18.} Keton jenuh dan keton $\alpha,\beta$-tak jenuh (sebuah sikloheksenon).
\textbf{19.} Karbon fusi cincin yang membawa gugus metil: empat substituen yang berbeda.
\textbf{20.} Tidak: dionnya akiral, dan kedua karbonilnya diserang secara
sama (keduanya saling merupakan bayangan cermin terhadap rantai); dengan
pereaksi akiral, produknya rasemat. Katalis amina kiral lebih menyukai satu
karbonil dan menghasilkan satu enansiomer secara berlebih (jilid Tahun 3).
\textbf{21.} Dion \ce{C7H10O2}: \qty{126.155}{g/mol}; \ce{C4H6O}: \qty{70.091}{g/mol}; \ce{C11H14O2}:
\qty{178.231}{g/mol}.
\textbf{22.} \ce{C7H10O2 + C4H6O -> C11H14O2 + H2O}: C 11, H 16, O 3 di setiap ruas.
\textbf{23.} $10.0/126.155 = \qty{0.07927}{mol} = \qty{79.3}{mmol}$.
\textbf{24.} $0.07927 \times 178.231 \times 0.70 = \boldsymbol{\approx \qty{9.89}{g}}$
keton Wieland--Miescher.
\end{solution}
